% About this book (front matter, unnumbered)
% Write in plain LaTeX. Code of a guide step: \guide{<step>}; commands: \begin{shell}...\end{shell};
% a gap in the draft: \placeholder{...}. See book/latex/queen.sty for the full list.

\chapter*{About this book}
\addcontentsline{toc}{chapter}{About this book}
\markboth{About this book}{About this book}

This book is written by the authors of QueenMQ. In it, we have tried to explain in simple terms
what QueenMQ does, and why. Here is what each chapter covers.

\paragraph{1. Introduction.} What QueenMQ is, why we built it (the story that went from a Postgres
table to a log on disk replicated with Raft), and the model behind it, in four ideas: partitions,
leases, transactions and a simple interface.

\paragraph{2. Writing applications with QueenMQ.} Every feature of the client API, one at a time, in
JavaScript: a local broker, queues and their options, push, pop and consume, acks, retries and the
dead-letter queue, transactions, KV, locks and timers, ephemeral queues and stream processing. Each
feature comes with the reason it exists.

\paragraph{3. Real applications examples.} The common messaging patterns, one page each: work queue,
publish/subscribe, request and reply, ordering and batching, workflows, windowed aggregation, rate
limiting, last value, delayed messages, retries and dead letters, and one worker at a time. Each
one solves a real problem in a few lines.

\paragraph{4. Integrations.} How QueenMQ works with the systems around it: Kafka clients through
its built-in Kafka port, an S3 bucket kept up to date by the sink, and PostgreSQL tables read and
written by the source and the sink.

\paragraph{5. Manage clusters.} The settings that matter most, a cluster with Docker and inside
Kubernetes, the configuration for different needs (tenants included), upgrades and monitoring, and
what to do when something breaks.

\paragraph{6. Architecture.} How the broker works inside: one leader per tenant, the path of a write
from the request to the disk, Raft, the memory state, the consumption engine, LMDB, and the files
on disk down to their bytes.

\medskip
Every piece of code in this book comes from a program that runs it against QueenMQ 2.6.0 and checks
the result. The examples are in JavaScript; the website has them in every language of the SDKs.

We hope this book helps everyone make the most of QueenMQ.
